% !PRINTED 611
\weberpage{611}

%<<CONT
integral linear substitution with the determinant $\pm 1$ (\S.\,147), and consequently the two forms

\[ (r\omega+s)\,ec\,\varphi_1,\quad e_1c_1\varphi \]

are associated with one another, and this has the consequence that the two forms $\varphi,\varphi_1$ are equivalent in the sense defined in \S.\,152. With this we have proved:

2. \gesperrt{If the quadratic irrational numbers $\omega,\omega_1$ are equivalent, then the corresponding functionals $\varphi,\varphi_1$ are also equivalent.}

But the converse theorem can also be proved:

3. \gesperrt{If the functionals $\varphi,\varphi_1$ are equivalent, then the quadratic irrational numbers $\omega$ and $\omega_1$ are also equivalent.}

To carry out this proof, we denote by $\psi$ a functional of the class that is reciprocal to the class of $\varphi$, so that $\varphi\psi$ is equivalent to a number. Then we can set:

\weq{5}{\varphi\psi = \varepsilon\xi,}

in which $\varepsilon$ is a unit and $\xi$ a \gesperrt{number} in $\mathfrak{o}$. Herein we can assume $\varphi$ and $\psi$ both to be basis forms, and accordingly set

\weq{6}{\varphi = ec(t_1+t_2\omega),\quad \psi = e'c'(t_1+t_2\eta)}

in which $e', c', \eta$ have the same meaning for $\psi$ as $e,c,\omega$ have for $\varphi$. For abbreviation we set [\S.\,166, (13)]:

\weq{7}{N_a(\varphi) = e^2c = P,\quad N_a(\psi) = e'^2c' = Q.}

If we denote by $\psi'$ the form conjugate to $\psi$, and for the moment set the unit form

\[ c't_1^2 + b't_1t_2 - a't_2^2 = \varepsilon_1, \]

then it follows from (5) by multiplication with $\psi'$:

\[ Q\varepsilon_1\varphi = \varepsilon\xi\psi', \]

and consequently

\weq{8}{\frac{Q\varphi}{\xi} = \frac{\varepsilon}{\varepsilon_1}\psi' = \chi,}

in which $\chi$ and $\psi'$ are associated. If we set

\weq{9}{\frac{Qec}{\xi} = \beta_1,\quad \frac{Qec\omega}{\xi} = \beta_2,}
%CONT>>
